\clearpage
\section{Testing Strategy Encyclopedia}

Testing strategy is where the abstract quality claims of a UI library are translated into repeatable engineering practice. In mature teams, testing is not a postscript to implementation; it is the mechanism that preserves component contracts as the design system evolves.

\subsection{Testing philosophy}
A good UI testing strategy separates concerns into four layers: behavior, accessibility, visual fidelity, and operational resilience. Unit tests should prove component logic. Integration tests should prove composition. Visual tests should prove appearance. End-to-end tests should prove workflows. None of these replaces the others.

The highest-quality teams also treat testing as a design-system property. A component that cannot be tested predictably is not truly reusable, because every reuse increases uncertainty.

\subsection{Unit testing component behavior in each library}
Unit tests should focus on isolated behavior: toggles, validation state, disabled state, keyboard interaction, and derived output. The exact technique differs by library, but the goal is consistent.

\subsubsection{shadcn/ui and local primitives}
Source-owned primitives are easy to test because the component source is local and explicit. Teams usually test render state, event callbacks, and variant output.

\begin{verbatim}
import { render, screen } from "@testing-library/react"
import userEvent from "@testing-library/user-event"
import { Button } from "@/components/ui/button"

test("calls onClick once", async () => {
  const user = userEvent.setup()
  const onClick = vi.fn()

  render(<Button onClick={onClick}>Save</Button>)
  await user.click(screen.getByRole("button", { name: "Save" }))

  expect(onClick).toHaveBeenCalledTimes(1)
})
\end{verbatim}

\subsubsection{Material UI}
MUI unit tests should check prop-driven rendering and theme-dependent behavior. Because MUI often uses wrappers, regression risk often appears in custom theme overrides.

\begin{verbatim}
import { ThemeProvider, createTheme } from "@mui/material/styles"
import { render, screen } from "@testing-library/react"

const theme = createTheme({ palette: { primary: { main: "#0f172a" } } })

render(
  <ThemeProvider theme={theme}>
    <MyButton />
  </ThemeProvider>
)
\end{verbatim}

\subsubsection{Chakra UI}
Chakra unit tests frequently need a provider wrapper because tokens, color mode, and system context influence rendering.

\begin{verbatim}
function renderWithChakra(ui: React.ReactElement) {
  return render(<ChakraProvider>{ui}</ChakraProvider>)
}
\end{verbatim}

\subsubsection{Ant Design}
Ant Design tests should check form state, table behavior, and controlled component outputs. Enterprise workflows often fail through subtle interaction regressions rather than pure rendering defects.

\begin{verbatim}
test("shows validation message", async () => {
  render(<LoginForm />)
  await userEvent.click(screen.getByRole("button", { name: /sign in/i }))
  expect(await screen.findByText(/email is required/i)).toBeInTheDocument()
})
\end{verbatim}

\subsubsection{Radix UI and custom systems}
Radix primitives are often tested by asserting behavior contracts: open/close state, keyboard navigation, and ARIA exposure.

\begin{verbatim}
expect(screen.getByRole("dialog")).toHaveAttribute("aria-modal", "true")
\end{verbatim}

\subsection{Integration testing with React Testing Library}
React Testing Library is the default integration layer for most React teams because it encourages user-centered assertions.

\subsubsection{Testing the rendered experience, not the implementation}
The best practice is to query by role, label text, and visible output. Avoid brittle selectors unless the UI surface is genuinely non-semantic.

\begin{verbatim}
const user = userEvent.setup()
render(<SearchBox />)
await user.type(screen.getByRole("searchbox"), "invoice")
expect(screen.getByText(/results/i)).toBeVisible()
\end{verbatim}

\subsubsection{Composition tests for shared shells}
Integration tests should verify that sidebars, headers, drawers, modals, and table filters compose correctly.

\begin{verbatim}
render(<DashboardLayout />)
expect(screen.getByRole("navigation")).toBeInTheDocument()
expect(screen.getByRole("main")).toBeInTheDocument()
\end{verbatim}

\subsection{Accessibility testing with axe-core automation}
Axe automation should run in component tests and in CI for critical pages. The value is not just compliance; it is regression detection for the easiest-to-miss defects.

\begin{verbatim}
import { axe } from "jest-axe"

it("has no axe violations", async () => {
  const { container } = render(<LoginForm />)
  expect(await axe(container)).toHaveNoViolations()
})
\end{verbatim}

For libraries with headless primitives, axe tests should cover focus visibility, label association, and dialog semantics. For libraries with opinionated visuals, axe tests should also cover theme contrast and color-mode transitions.

\subsection{Visual regression testing with Chromatic and Percy}
Visual regression testing is essential when teams depend on design-system stability. Chromatic is especially useful when Storybook is the source of truth; Percy works well when teams need production-route screenshots and broader workflow coverage.

\begin{itemize}[leftmargin=*, itemsep=0.35em]
  \item Use Chromatic for component-level visual diffs in Storybook.
  \item Use Percy for cross-route and cross-browser visual baselines.
  \item Keep snapshot thresholds strict for foundational primitives.
  \item Allow slightly wider tolerances for content-rich and data-heavy surfaces.
\end{itemize}

\subsection{End-to-end testing with Playwright and Cypress}
E2E testing should verify user journeys that span authentication, navigation, filtering, and saving state. Playwright is strong for multi-browser automation; Cypress is strong for fast feedback and rich debugging in app-level workflows.

\begin{verbatim}
// Playwright example
await page.getByRole("button", { name: "Sign in" }).click()
await expect(page.getByRole("heading", { name: /dashboard/i })).toBeVisible()
\end{verbatim}

\begin{verbatim}
// Cypress example
cy.findByRole("button", { name: /save/i }).click()
cy.findByText(/saved successfully/i).should("be.visible")
\end{verbatim}

\subsection{Storybook-driven development workflow}
Storybook provides an excellent feedback loop when used as a design-system workbench rather than a demo gallery. Stories should encode state combinations, edge cases, and accessibility notes.

\begin{itemize}[leftmargin=*, itemsep=0.35em]
  \item Create one story per meaningful state.
  \item Add controls only when they clarify component APIs.
  \item Use MDX or docs blocks to explain conventions and pitfalls.
  \item Wire Chromatic to every merge request for visible regression gating.
\end{itemize}

\subsection{Testing custom hooks that interact with theme context}
Hooks that depend on theme values, color mode, or token configuration should be wrapped in a provider test harness.

\begin{verbatim}
function renderHookWithTheme<T>(hook: () => T) {
  return renderHook(hook, {
    wrapper: ({ children }) => <ThemeProvider>{children}</ThemeProvider>,
  })
}
\end{verbatim}

\subsection{Mocking library providers in tests}
Mock only the providers that are required to isolate the behavior under test. Over-mocking hides integration failures.

\begin{verbatim}
vi.mock("@/lib/auth", () => ({
  useSession: () => ({ user: { id: "1" }, loading: false }),
}))
\end{verbatim}

\subsection{Snapshot testing tradeoffs and when to avoid it}
Snapshots are useful for stable, low-entropy output but harmful when overused for dynamic layouts or content-rich screens. Avoid snapshots for tables with user-generated content, charts, or components with unavoidable date/time variance.

A better approach is to snapshot stable structure and assert dynamic behavior separately.

\subsection{Component contract testing for design systems}
Contract tests verify that shared components continue to honor the public API expected by application teams.

\begin{itemize}[leftmargin=*, itemsep=0.35em]
  \item Required props remain required.
  \item Default variants do not change silently.
  \item Keyboard semantics remain stable.
  \item Theme tokens continue to resolve correctly.
\end{itemize}

\subsection{Performance testing components under load}
Performance tests should stress the cases that break at scale: large tables, virtualized lists, long trees, and repeated rerenders.

\begin{verbatim}
for (let i = 0; i < 1000; i++) {
  rerender(<VirtualizedTable rows={bigDataset.slice(0, i)} />)
}
\end{verbatim}

Measure commit time, interaction latency, and memory growth over repeated mount/unmount cycles.

\subsection{Cross-browser testing strategy}
Cross-browser coverage should include Chromium, WebKit, and Firefox where possible. Focus on layout, font rendering, keyboard behavior, and scroll mechanics.

\begin{table}[htbp]
  \centering
  \caption{Cross-browser focus areas}
  \begin{tabularx}{\textwidth}{@{}p{0.24\textwidth}Y@{}}
    \toprule
    \textbf{Browser family} & \textbf{Primary risk area} \\
    \midrule
    Chromium & Fast-moving feature support and layout assumptions. \\
    Firefox & Scroll, focus, and CSS edge-case behavior. \\
    WebKit & Mobile rendering, input behavior, and color/zoom quirks. \\
    \bottomrule
  \end{tabularx}
\end{table}

\subsection{Mobile testing considerations}
Mobile testing should verify touch targets, gesture handling, virtual keyboard overlap, and scrolling behavior inside drawers or dialogs. A desktop-perfect component can still fail badly on small screens.

\subsection{Testing dark mode and theme switching}
Theme switching tests should assert token propagation, contrast stability, and persistence of user preference.

\begin{verbatim}
render(<ThemeToggle />)
await userEvent.click(screen.getByRole("switch"))
expect(document.documentElement).toHaveClass("dark")
\end{verbatim}

\subsection{Continuous integration setup for component testing pipelines}
A mature CI pipeline should run in layers: lint, typecheck, unit tests, accessibility checks, story snapshots, visual diffs, and E2E smoke flows.

\begin{enumerate}[leftmargin=*, itemsep=0.3em]
  \item Fast preflight: type and lint.
  \item Component-level confidence: unit, RTL, axe.
  \item Visual gate: Storybook and image diff.
  \item Workflow gate: Playwright or Cypress smoke tests.
  \item Release gate: browser matrix and performance budgets.
\end{enumerate}

\begin{highlightbox}{Testing principle}
The best testing strategy is layered, user-centered, and selective. It protects the design system without turning every component into a fragile mass of duplicated assertions.
\end{highlightbox}
